\par
\SourceLeaf{488}%
Ich habe noch hinzuzufügen, dass ich Sr. — — — über diesen Gegenstand bereits vor einem Jahr eine ausführliche Denkschrift vorgelegt habe. Er erkannte die Bedeutung der Fortsetzung der dänischen Messungen und sicherte wenigstens sofort die Möglichkeit dazu, indem er mir befahl, Lüneburg an die südlichen Punkte jener Messung anzuschließen, die zum Teil aus vergänglichen Signalpunkten bestanden. Dies ist im vergangenen Herbst geschehen und hat die künftige Fortsetzung vom Verlust der gefährdeten Punkte unabhängig gemacht.

Wenn Ew. — — diesen Ideen Ihren Beifall schenken und sie für wert halten, sich für ihre Verwirklichung einzusetzen, muss ich davon ausgehen, dass eine möglichst vollkommene Ausführung beabsichtigt ist, die neben jeder anderen ehrenvoll bestehen kann. In diesem Sinne würde ich die Ausführung, wenn sie mir anvertraut würde, gerne übernehmen und eine vielleicht länger als einen Sommer dauernde Unterbrechung meiner rein astronomischen Arbeiten für kein Opfer halten.

Was den Kostenaufwand für die Verlängerung der dänischen Gradmessung durch das Königreich Hannover betrifft, lässt er sich wegen seiner Abhängigkeit von mancherlei unvorhersehbaren Umständen freilich unmöglich mit einiger Genauigkeit im Voraus angeben. So hat die Witterung einen wesentlichen Einfluss auf die Dauer und damit auf die Kosten. Ich wüsste von keiner Gradmessung, deren Kosten öffentlich bekannt gemacht worden wären. Auch wenn ich vom dänischen Astronomen Auskunft über die bisherigen Kosten seiner Gradmessung erhalten könnte, ließen sich daraus die Kosten ähnlicher Unternehmungen im Königreich Hannover wegen der großen Verschiedenheit der örtlichen Verhältnisse doch nur sehr unsicher schätzen. So werden etwa in Hannover die Fuhrkosten für den Transport von Personen und Instrumenten einen sehr beträchtlichen Teil der Gesamtkosten ausmachen. In Dänemark fallen diese Kosten fast ganz aus der Rechnung, da die Fuhrleistungen dort ex officio in natura erbracht werden. Doch glaube ich, alles wohl erwogen, dass die Summe von 1500 Pfund Sterling zur Deckung aller Kosten ausreichen würde. Dass darüber anschließend Rechnung gelegt würde, versteht sich von selbst.

Professor \textsc{Schumacher} hat bei seiner Gradmessung außer ein paar Schreibgehilfen zwei Offiziere im Hauptmannsrang als Gehilfen. Ihre Aufgabe ist es, die Gegend vorher zu bereisen, geeignete Punkte für Dreiecksstationen auszusuchen und sie nach vorläufigen Messungen vor Ort Prof. S. zur Auswahl vorzuschlagen. Danach sorgen sie an den ausgewählten Punkten für die Errichtung von Signalen, wo es nötig ist, und treffen die übrigen notwendigen Vorkehrungen. Mit einem Wort, sie treffen alle Vorbereitungen, damit der Astronom überall ohne großen Zeitverlust mit der Beobachtung beginnen kann. Schließlich übernehmen sie, wo es nötig ist, bei den Beobachtungen die Nebenarbeiten. Für solche Gehilfen kommt es daher weniger auf besonders tiefe mathematische oder astronomische Einsichten an als vielmehr auf lebhaften Eifer für die Sache, größte Pünktlichkeit, Sinn für höchste Genauigkeit, eine gewisse praktische Geschicklichkeit, einige Kenntnisse im Bauwesen und eine gewisse Vertrautheit mit dem Geschäftsgang der Behörden in unserem Land, mit denen man in solchen Angelegenheiten zu tun hat. Auch für die hannoversche Gradmessung wäre kein größeres Personal nötig als für die dänische. Sobald ein Beschluss gefasst ist, würde ich mich unverzüglich darum bemühen, selbst geeignete Gehilfen auszusuchen. Natürlich läge es aber im Ermessen von Ew. — —, ob es vielleicht ratsam wäre, \textit{in der oben erwähnten Hinsicht} noch mehrere Personen als Freiwillige bei den Messungen zugegen sein zu lassen, damit sie sich auf anspruchsvollere geodätische Arbeiten vorbereiten können. — —

Göttingen, den 30. Mai 1819.
\begin{flushright}\textsc{Carl Friedrich Gauss}.\end{flushright}
